% Generated draft: requires mathematical and source review.
\documentclass[11pt]{article}
\usepackage[margin=25mm]{geometry}
\usepackage{fontspec}
\setmainfont{Latin Modern Roman}
\usepackage{amsmath,amssymb,mathtools,unicode-math}
\setmathfont{Latin Modern Math}
\usepackage{enumitem,needspace,xurl,hyperref}
\setlength{\parindent}{0pt}
\setlength{\parskip}{5pt}
\setlength{\emergencystretch}{3em}
\newcommand{\N}{\mathbb{N}}
\newcommand{\Z}{\mathbb{Z}}
\newcommand{\Q}{\mathbb{Q}}
\newcommand{\R}{\mathbb{R}}
\newcommand{\C}{\mathbb{C}}
\newcommand{\F}{\mathbb{F}}
\newcommand{\A}{\mathbb{A}}
\newcommand{\PP}{\mathbb{P}}
\newcommand{\E}{\mathbb{E}}
\newcommand{\eps}{\varepsilon}
\newcommand{\dd}{\,\mathrm{d}}
\begin{document}
\section*{UnsolvedMath ID: 30001589}
\textbf{Source problem code:} catalogue entry
\subsection{Definitions and mathematical statement}
Let \(X\) be a geodesic metric space that is \(\operatorname{CAT}(0)\): distances between points on a geodesic triangle are no greater than the corresponding distances in its Euclidean comparison triangle. Let \(G\) act on \(X\). The action is free if \(gx=x\) implies \(g=1\), and cocompact if the quotient \(G\backslash X\) is compact. An isometric Euclidean plane in \(X\) is a map \(\iota:\R^2\to X\) satisfying
\[
 d_X(\iota(u),\iota(v))=\|u-v\|_2
 \qquad\text{for all }u,v\in\R^2.
\]
A copy of \(\Z^2\) in \(G\) means a subgroup isomorphic to the additive group \(\Z^2\), equivalently an injective group homomorphism \(\Z^2\to G\).

\textbf{Flat closing question, subject to an unspecified class.} For every group and action belonging to the source's intended class of ``random intermediate-rank groups'' acting freely and cocompactly on \(\operatorname{CAT}(0)\) spaces, is the implication
\[
 \bigl(\exists\text{ an isometric embedding }\iota:\R^2\hookrightarrow X\bigr)
 \quad\Longrightarrow\quad
 \bigl(\exists\text{ an injective homomorphism }\varphi:\Z^2\hookrightarrow G\bigr)
\]
valid? This displays the deterministic implication being asked about; the supplied record does not determine whether its assertion is intended for every realization, almost surely, or in some other probabilistic sense.

\textbf{Further questions.} For groups in the same intended class, what rigidity properties hold? Are these groups linear, and are they residually finite? In the usual unrestricted sense, linearity means that there exist a field \(K\), an integer \(n\geq 1\), and an injective homomorphism \(G\to\operatorname{GL}_n(K)\). Residual finiteness means that for every \(g\in G\setminus\{1\}\), there exist a finite group \(F\) and a homomorphism \(q:G\to F\) with \(q(g)\neq 1\).

\textbf{Unresolved qualifications.} The supplied record defines neither intermediate rank nor the random construction, its parameters, or its probability law. It does not explicitly require the action to be by isometries or to be proper, or require \(X\) to be proper or complete; these are not added as hypotheses here. It specifies neither a rigidity property nor a comparison class, and gives no field or characteristic restriction for linearity. Consequently, the displayed implication is a conditional formulation of the supplied question, not a fully specified assertion about a defined family of groups. The further questions also require clarification of their class and quantifiers.
\subsection{Short English statement}
The source asks whether an isometric Euclidean plane in the space of a free, cocompact action of a random intermediate-rank group forces that group to contain \(\Z^2\). It also asks about rigidity, linearity, and residual finiteness. The group family, probabilistic quantifiers, and intended rigidity properties are unspecified, so this remains a statement-only draft requiring mathematical review.
\subsection{Sources}
\begin{itemize}
\item Local UnsolvedMath registry: permanent ID \(30001589\), source problem code \texttt{catalogue entry}, title ``Flat Closing and Rigidity for Intermediate-Rank Groups,'' in the Oberwolfach Reports Open Problems collection. No individual proposer is supplied; the proposed year is recorded as 2010.
\item Supplied source citation: \emph{Actions and Invariants of Residually Finite Groups: Asymptotic Methods} (2011), \emph{Oberwolfach Reports}, pp.~2335--2391, DOI \texttt{10.4171/owr/2010/40}. \url{https://doi.org/10.4171/owr/2010/40}. The record describes Oberwolfach Reports as documenting mathematical workshops and their open-problem discussions. It gives no problem-specific page or named contributor.
\item Supplied literature source: \url{https://arxiv.org/abs/1112.4637}. No title, authors, or specific result from this reference are supplied.
\item Supplied literature assessment dated \texttt{2026-08-21}: the record labels the problem open and states that the flat closing question remains open in the general \(\operatorname{CAT}(0)\) group form, with results under partial hypotheses and in special cases. This is supplied dated curation derived from the OWR record, not an independently verified literature assessment; its time-sensitive status and cited claims require independent verification.
\item Supplied statement-verification metadata: \texttt{corrected\_verified}. The record reports that the mathematical content was checked against the official Oberwolfach report and that notation, context, and question form were normalized without changing the source-level claim. This verification report is supplied metadata, not verification performed for this draft.
\end{itemize}
\end{document}
